% Source: https://artofproblemsolving.com/community/c3h2501202_is_0_negative_positive_or_neither
% Stable identity: aops-topic-2501202-q01
\textbf{(a)} Whether a number is added or subtracted says nothing about its sign: positive and negative numbers can both be added and subtracted. For example, in $5+(-3)$ a negative number is added, and in $12-3$ a positive number is subtracted. Subtracting a number means adding its opposite, $a-b=a+(-b)$, so $12-0=12+(-0)=12+0$, since $-0=0$. The sign of a number is a property of the number itself, decided by comparing it with $0$, not by the operation in which it appears.

\textbf{(b)} For every real number $x$, exactly one of $x<0$, $x=0$, $x>0$ is true. For $x=0$ the middle statement is true, so $0>0$ and $0<0$ are both false: $0$ is neither positive nor negative. If $x\ne0$, then $x>0$ or $x<0$, so $x$ is positive or negative. Hence $0$ is the only real number that is neither.

\textbf{(c)} If $x=-x$, then $2x=0$, so $x=0$; and indeed $-0=0$. If $x>0$, adding $-x$ to both sides gives $0>-x$, so $-x$ is negative. Now suppose that $0$ were positive. Then $-0$ would be negative; but $-0=0$, so $0$ would be both positive and negative, which is impossible by (b). So $0$ is not positive (and, in the same way, using that $-x$ is positive when $x$ is negative, $0$ is not negative). $\blacksquare$
